% RESEARCH_DRAFT_SHA256: 3232fedfd7aa3611e05219242a2e3fd484d62701e56e1103bf4cccd7ebbad925
% Layout and source-reference wrappers expanded for the web renderer.
\documentclass[11pt]{article}
\usepackage[margin=25mm]{geometry}
\usepackage{fontspec,unicode-math}
\setmainfont{Latin Modern Roman}
\setmathfont{Latin Modern Math}
\usepackage{amsmath,amssymb,mathtools,xurl,hyperref}
\hypersetup{colorlinks=true,urlcolor=blue,linkcolor=blue}
\setlength{\parindent}{0pt}
\setlength{\parskip}{5pt}
\setlength{\emergencystretch}{3em}
\newcommand{\sref}[2]{\hyperlink{source:#1:#2}{[S#2]}}
\newcommand{\cataloguescope}[1]{\paragraph{Source question.}#1}
\newcommand{\floor}[1]{\left\lfloor#1\right\rfloor}
\begin{document}

\section*{Adaptive source coding: countable mixtures and the cost of uniform constants}
\noindent UnsolvedMath ID 30003172; OWR-14741-004. Research attempt dated 11 October 2026.
\paragraph{Result in scope.}This manuscript contains a source-specific attempt and complete proofs of the stated auxiliary calculations. It makes no new solution or novelty claim for the entire target.
\section{Original problem and preserved definition}
The following definition and source list are reproduced from the repository. Any earlier inspection dates and reported literature claims in this preserved text describe that source record, not additional verification by this attempt.
% BEGIN PRESERVED DEFINITION
\subsection{Definitions and mathematical statement}

\paragraph{Shared notation.}Standard number systems are denoted by $\mathbb N,\mathbb Z,\mathbb Q,\mathbb R,\mathbb C$; index ranges, computational models and specialized notation follow the cited problem source and the definitions below.


A coding probability $Q_n$ assigns probabilities to length-$n$ source blocks. Redundancy is the excess expected codelength relative to the source entropy; minimax redundancy minimizes the worst such excess over a class.

A code $C\subseteq\Sigma^N$ is a set of permitted length-$N$ words. Its Hamming minimum distance is the minimum number of unequal coordinates of two distinct words. For a linear $[N,K,D]_q$ code, rate is $K/N$ and relative distance is $D/N$; decoding recovers a codeword from corrupted observations.

For a discrete random variable $X$, entropy is $H(X)=-\sum_xp(x)\log p(x)$, and mutual information is $I(X;Y)=H(X)+H(Y)-H(X,Y)$. Coding rates, capacities, and exponents depend on the stated channel or source model and error criterion.
\paragraph{Statement recorded in the supplied source.}
Let $(C(\alpha))_{\alpha\in A}$ be source subclasses with nontrivial minimax redundancy rates. Does there exist a single sequence of coding probabilities $Q_n$, independent of $\alpha$, whose redundancy is asymptotically optimal, or optimal within a fixed factor, simultaneously for every subclass in the collection? \href{https://doi.org/10.4171/owr/2016/33}{[S2]}
\paragraph{Scope qualification.}
The source collection, redundancy convention, and exact constant-factor uniformity are the report\textquotesingle s assumptions.
\paragraph{Source status.}
The supplied catalogue labels this record partially solved. Its source-status metadata was checked against the supplied snapshot on 2026-10-11. This import does not independently establish current open status or mathematical correctness. The catalogue literature assessment is dated 2026-08-21.
\subsection{Short English statement}
Find one code that attains each source subclass' s optimal redundancy rate without knowing the subclass.
\subsection{Sources}
\begin{itemize}
\item[S1] Ulam AI and the UnsolvedMath Contributors, UnsolvedMath: A Curated Collection of Open Mathematics Problems, record 30003172 (OWR-14741-004). Catalogue curation and metadata: CC BY 4.0. The supplied source question is retained; definitions and the English summary are editorial additions.
\url{https://huggingface.co/datasets/ulamai/UnsolvedMath}.
\item[S2] Learning Theory and Approximation (2017), Oberwolfach Reports, DOI 10.4171/owr/2016/33, pp. 1875--1941
\url{https://doi.org/10.4171/owr/2016/33}.
\item[S3] Ben-Hamou--Boucheron--Gassiat, Pattern coding meets censoring
\url{https://arxiv.org/abs/1608.08367}.
\item[S4] Boucheron--Gassiat--Ohannessian, Adaptive coding
\url{https://arxiv.org/abs/1402.6305}.
\textit{Underlying Oberwolfach Reports material retains its source-specific licence, including CC BY-SA 4.0 where applicable; see THIRD\_PARTY\_NOTICES.md.}
\end{itemize}
% END PRESERVED DEFINITION
\subsection{Verbatim supplied catalogue question}
Let $(C(\alpha))_{\alpha\in A}$ be source subclasses with nontrivial minimax redundancy rates. Does there exist a single sequence of coding probabilities $Q_n$, independent of $\alpha$, whose redundancy is asymptotically optimal, or optimal within a fixed factor, simultaneously for every subclass in the collection?

\section{The source convention matters}
Boucheron's report, printed pages 1921--1924, concerns stationary
memoryless sources on a countable alphabet, and uses block redundancy
$D(P^n\Vert Q_n)$, not channel-code error or Hamming distance.
It allows a constant approximation factor depending on the subclass
index. Its original theorem gives a $\log\log n$ factor for regularly
varying envelope classes; it does not state full constant-factor
adaptivity over every such envelope. We use natural logarithms below,
so changing the log base scales every redundancy by the same constant.
The mixture route and bounded-redundancy example were suggested by
the coordinating Codex GPT-6 model and independently checked within
this model-assisted attempt. They are standard known mechanisms,
not claimed novel source-coding theorems.

\section{A full countable-mixture sufficiency theorem}
Let $C_j$, $j\ge1$, be a countable collection of source classes and
let $Q_{j,n}$ be coding probabilities for each block length. Choose
fixed $w_j>0$ with $\sum_jw_j=1$ and put
\[
Q_n(x)=\sum_{j\ge1}w_jQ_{j,n}(x).
\]
This is a probability distribution: sums over the countable block
alphabet can be interchanged for nonnegative terms, giving total mass
$\sum_jw_j=1$. For every $j$, $Q_n\ge w_jQ_{j,n}$ pointwise.
Thus for any source $P$,
\[
D(P^n\Vert Q_n)\le D(P^n\Vert Q_{j,n})+\log(1/w_j).
\]
To justify the inequality even on an infinite alphabet, when the
right-hand divergence is infinite it is trivial; otherwise integrate
the pointwise logarithmic ratio inequality. Negative parts are
integrable: on entries with $P^n(x)<Q_n(x)$, writing
$z=P^n(x)/Q_n(x)\in[0,1]$ gives
$-P^n(x)\log z=Q_n(x)(-z\log z)\le Q_n(x)/e$,
and summing bounds their total by $1/e$. The same observation applies
to $Q_{j,n}$, so extended expectations are well defined and monotone
integration is legitimate. A zero coding mass at a positive source
mass gives the usual infinite divergence.

Define
\[
R_j^*(n)=\inf_Q\sup_{P\in C_j}D(P^n\Vert Q),
\qquad R(Q,C_j)=\sup_{P\in C_j}D(P^n\Vert Q).
\]
Taking suprema gives
$R(Q_n,C_j)\le R(Q_{j,n},C_j)+\log(1/w_j)$.
Suppose $0<R_j^*(n)<\infty$ eventually,
$R_j^*(n)\to\infty$, and
$R(Q_{j,n},C_j)\le(1+o_j(1))R_j^*(n)$ for each fixed $j$.
Divide the mixture inequality by $R_j^*(n)$; the fixed overhead
$\log(1/w_j)$ vanishes relatively. Since every coding probability
has worst-case redundancy at least the infimum $R_j^*(n)$, we conclude
\[
\frac{R(Q_n,C_j)}{R_j^*(n)}\longrightarrow1
\quad\text{for every fixed }j.
\]
The single sequence $Q_n$ is independent of $j$. The convergence need
not be uniform in $j$, and evaluating an infinite mixture need not be
computationally effective. The statement establishes existence under
its hypotheses, not an efficient online code for the entire source family.

\section{Bounded redundancy changes the conclusion}
On alphabet $\mathbb N$, let $C_j$ consist of the two iid point-mass
sources concentrated on symbols $2j$ and $2j+1$. Their length-$n$
outputs are the two distinct constant words $(2j)^n,(2j+1)^n$.
For every $n$ the minimax redundancy is exactly $\log2$.
Indeed assigning each word mass $1/2$ achieves that risk; conversely
their two masses sum to at most one, so one is at most $1/2$ and the
corresponding point-mass divergence is at least $\log2$.

There is no asymptotically optimal coding sequence for even two
distinct classes in this family. If there were, for each of the four
constant words it would have
$-\log Q_n(x)\le(1+o(1))\log2$, hence
$Q_n(x)\ge\tfrac12(1-o(1))$. Their sum would tend to at least two,
contradicting total mass one. There are only two classes in this
argument, so their separate asymptotic thresholds can be maximized;
no unjustified uniform limit is used.

More generally, a common factor $C<\infty$ for $J$ disjoint classes
would force each of their $2J$ words to have mass at least $2^{-C}$,
so $1\ge2J\,2^{-C}$ and $C\ge\log_2(2J)$.
For infinitely many classes, no one finite common factor can hold
eventually for every class. Given a proposed $C$, choose a finite
$J>2^{C-1}$ and then a block length beyond the largest of those
$J$ individual eventual thresholds to get the same contradiction.
If the bound is stated with an additional vanishing error, use
$C+1$ first and choose $J>2^C$.

The source, however, permits \emph{class-dependent} factors. They
do exist here: use the mixture of the two-word uniform distributions
with weights $w_j=2^{-j}$, $j\ge1$. Each word in class $j$ then has
mass $2^{-j-1}$, giving risk $(j+1)\log2$, exactly a factor $j+1$
of its minimax risk. Thus the example does not refute the source's
class-dependent constant-factor possibility. Nor does its bounded
redundancy refute the divergent-rate sufficiency theorem above.

\section{Where the mixture attempt stops for the envelope problem}
The source's family of envelope functions is not presented as a
countable list of classes. A probability distribution has positive
atoms at only countably many indices: for every integer $r\ge1$
there are at most $r$ atoms of mass at least $1/r$, and every positive
atom lies in one of these finite sets. Consequently assigning a
fixed positive individual weight to every member of an uncountable
family is impossible. A continuous mixture or countable approximation
could still work, but the pointwise inequality
$Q_n\ge w_jQ_{j,n}$ with fixed $w_j>0$ would have to be replaced
by a quantified approximation argument. In particular one would need
nearby dominating envelopes and control of their redundancy loss and
description overhead. Such an argument is not supplied here.

\section{Model-assisted audit and remaining scope}
The audit checked zero coding masses, infinite sums, logarithm
normalization, and the order of the fixed-index limit. Most crucially,
it separated a common constant from the class-dependent constants
explicitly allowed on source page 1922. The primary Pattern Coding
Meets Censoring abstract was checked against the report's stated
$\log\log n$ result; the full coding proof was not audited.
No numerical experiment or theorem prover was used. The countable
divergent-rate theorem and the bounded-rate obstruction are auxiliary;
the full collection and its efficient exact or constant-factor adaptive
code are not resolved by this attempt.

\section{Process and contributions}
Mathematics and model-assisted adversarial self-audit: OpenAI gpt-6-astra, ultra reasoning, as dispatched for this research batch; exact serving revision not exposed. Codex (GPT-6; exact serving revision not exposed) materially assisted with orchestration, format review, and the countable-mixture and bounded-redundancy mathematical suggestions described above. No human mathematical authorship, mathematical edits, or independent human proof checking is asserted.
The coordinating Codex (GPT-6; exact serving revision not exposed) also provided a separate model-assisted written-proof review of this manuscript. This is not independent human verification.
Contribution prompt A was applied to the full definition and supplied source row, then prompt B was applied adversarially to the derived results. The item-specific audit above records the scope checks and rejected extrapolations. Local assembly and structural checks do not establish mathematical correctness. The source catalogue status is preserved in TOP\_PROBLEM and is separate from this attempt status.
\paragraph{Integration review.} Codex (GPT-6; exact serving revision not exposed) reviewed the written arguments and expanded layout and source-reference wrappers. Repository metadata and the production web renderer were checked using Python 3.13 and JavaScript with MathJax 3.2.2. This model-assisted review is not independent human proof verification; no PDF compilation or formal proof checking is asserted.

\section{Completion Estimate}
\noindent\textbf{COMPLETION: 0\%}
The proved material is a rederivation of known or auxiliary facts rather than a substantiated new advance on the original target. The mixture argument proves countable divergent-rate adaptivity under supplied per-class coders, while bounded-rate examples separate exact adaptivity and uniform constants. It does not cover the full uncountable envelope collection or construct the requested efficient general adaptive code.
\section{References}
The complete original Sources subsection above is retained, including catalogue attribution and licence notices. The additional list records primary sources actually consulted for this attempt on 11 October 2026; uncited source links reproduced above were not all rechecked.
\begin{itemize}
\item Boucheron, Adaptive compression against countable alphabets, OWR 33/2016, pp. 1921--1924; full source definitions, factor convention and stated envelope theorem inspected in the cached primary report. \url{https://doi.org/10.4171/owr/2016/33}.
\item Ben-Hamou, Boucheron and Gassiat, Pattern Coding Meets Censoring (2016); primary abstract inspected, full coding proof not audited. \url{https://arxiv.org/abs/1608.08367}.
\end{itemize}
\section{Final status}
\textbf{Unresolved in this attempt.} The mixture argument proves countable divergent-rate adaptivity under supplied per-class coders, while bounded-rate examples separate exact adaptivity and uniform constants. It does not cover the full uncountable envelope collection or construct the requested efficient general adaptive code.
\end{document}
